This chapter addresses the work proposal for the dissertation. It proposes a testbed to establish an open, programmable environment that simplifies access to the network services offered by \acs{5G} systems, through standardized network exposure interfaces. This chapter identifies the overall design and the capabilities such an environment must provide so verticals can implement and validate their use cases. These are: (i) \acs{5G} resource lifecycle management, providing mechanisms to create, configure, and manage network slices and their associated \acsp{UE} to meet the heterogeneous requirements of each vertical (Section~\ref{chap:proposal:resource_lifecycle}); (ii) network exposure, letting verticals interact with the \acs{5G} control plane to request and configure the network according to their service requirements (Section~\ref{chap:proposal:exposure}); (iii) network observability, providing applications with relevant network data to inform decision-making and trigger appropriate actions (Section~\ref{chap:proposal:observability}); and (iv) service delivery, through which network slices and any services supporting the experiment can be ordered and managed on demand by third parties (Section~\ref{chap:proposal:delivery}). Integrating all these capabilities in a single end-to-end platform, whose architecture is defined in Section~\ref{chap:proposal:architecture}, advances the existing literature on future network testbeds, supporting interoperable service management aligned with \acs{TMF} \acs{ODA}. Section~\ref{chap:proposal:use-case} introduces a multi-vertical use case to demonstrate the interfaces and capabilities available in the testbed. Section~\ref{chap:proposal:capabilities} summarizes these capabilities, describing how the provided mechanisms support the requirements of the considered use case. Finally, Section~\ref{chap:proposal:validation} defines the validation methodology proposed to assess the resulting vertical application and, consequently, the exposed testbed capabilities.
